
\begin{frame}
\frametitle{MongoDB, système distribué \textit{maître-esclave}}


\figSlide{mongodb-ms}

\begin{redbullet}
  \item \mongodb/ gère la \alert{replication asynchrone} au sein d'un \textit{replica set} (RS)
  
  \item Le RS comprend un \textit{maître} (\textit{primary}) et des \textit{esclaves} (\textit{secondary})
  
   \item Si nombre pair\,: ajout d'un processus supplémentaire (\textit{arbiter})
   
   \item Les écritures ont \alert{toujours} lieu sur le \textit{maître}
   
   \item Le \textit{driver} associé à l'application peut lire sur le \textit{maître} (cohérence forte)
       ou sur un \textit{secondary} (cohérence faible, mais répartition)
  
\end{redbullet}

\end{frame}


\begin{frame}
\frametitle{Election d'un \textit{maître} dans \mongodb/}

Une élection est déclenchée dans les cas suivants\,:

\begin{redbullet}
  \item Initialisation d'un nouveau RS
  \item Un \textit{esclave} perd le contact avec le \textit{maître}
  \item Le \textit{maître} perd son statut car il ne voit plus qu'une minorité de participants.
\end{redbullet}
\vfill

\begin{redbullet}
  \item Tous les clients (\textit{drivers}) connectés au RS sont réinitialisés pour dialoguer avec le nouveau \textit{maître}.
  \item Beaucoup d'options pour paramétrer le vote (p.e., une \alert{priorité} donnée à chaque serveur). 
 \end{redbullet}

\end{frame}


\begin{frame}
\frametitle{Cohérence et répartition des lectures}

Par défaut, les clients lisent sur le \textit{maître}\,: \alert{cohérence forte}.

\vfill

Une option de connexion permet à l'application d'indiquer qu'elle accepte de lire sur un \textit{esclave}\,:
   \alert{cohérence faible}
   
\vfill
   
Autre option\,: l'application peut attendre que $n$ écritures soient faites avant de reprendre la main.

\begin{remblock}{Remarque}
La réplication n'est pas, dans \mongodb/, le moyen privilégié de passer à l'échelle. Le \alert{partitionnement}
est un mécanisme plus puissant (à suivre!)
\end{remblock}

\end{frame}


\begin{frame}[fragile]
\frametitle{A l'action\,! Test de la réplication}

On crée trois conteneurs avec des serveurs MongoDB

\begin{minted}{text}
docker run --name mongo1 --net host mongo mongod --replSet mon-rs --port 30001
docker run --name mongo2 --net host mongo mongod --replSet mon-rs --port 30002
docker run --name mongo3 --net host mongo mongod --replSet mon-rs --port 30003
\end{minted}
\vfill

Connectons un client au premier serveur (ou Studio3D si vous voulez).


\begin{minted}{text}
mongo --host 192.168.99.100 --port 30001
\end{minted}

Et initialisons le RS 

\begin{minted}{text}
rs.initiate()
\end{minted}

On peut alors ajouter le second n{\oe}ud (remplacer l'IP de la machine).

\begin{minted}{text}
rs.add ("192.168.99.100:30002")
\end{minted}


\end{frame}

\begin{frame}[fragile]
\frametitle{Continuons}

La fonction suivante indique le statut du n{\oe}ud auquel on est connecté.

\begin{minted}{text}
db.isMaster()
\end{minted}

Information plus complète sur le RS avec:

\begin{minted}{text}
rs.status()
\end{minted}

Tout va bien\,? Alors insérons.
\begin{minted}{text}
use nfe204;
db.test.insert ({"produit": "Grulband", prix: 230, enStock: true})
\end{minted}

On devrait trouver le document sur chaque \textit{esclave}\,?
\begin{minted}{text}
mongo --host 192.168.99.100 --port 30002
use nfe204;
db.test.find ()
\end{minted}

Hmmm que se passe-t-il... \,? Et après \texttt{rs.slaveOk()}\,? Pourquoi donc\,?

\end{frame}


\begin{frame}[fragile]
\frametitle{Et la reprise sur panne\,?}

Tuons (gentiment) notre \textit{maître}.

\begin{minted}{text}
db.shutdownServer()
\end{minted}

Question\,: le second va-t-il s'élire \textit{maître}\,?

\vfill

Relancer le premier serveur. Que se passe-t-il\,?
\vfill

\begin{remblock}{Remarque}
Un bon \textit{driver} devrait enregistrer la liste des serveurs du RS, et se reconnecter
automatiquement au \textit{maître} quand celui-ci change.
\end{remblock}


\end{frame}

\end{document}


